\documentclass[12pt,a4paper]{article}

% ─── Pacotes ───────────────────────────────────────────────────────────────────
\usepackage[brazil]{babel}
\usepackage[utf8]{inputenc}
\usepackage[T1]{fontenc}
\usepackage{geometry}
\usepackage{amsmath, amssymb}
\usepackage{booktabs}
\usepackage{array}
\usepackage{parskip}
\usepackage{indentfirst}
\usepackage{setspace}
\usepackage{titlesec}
\usepackage{algorithm}
\usepackage{algpseudocode}
\usepackage{hyperref}
\usepackage{xcolor}

\geometry{
  a4paper,
  top=3cm, bottom=2cm,
  left=3cm, right=2cm
}

\onehalfspacing

\hypersetup{
  colorlinks=true,
  linkcolor=black,
  urlcolor=blue
}

% ─── Cabeçalho ─────────────────────────────────────────────────────────────────
\begin{document}

\begin{center}
  {\large \textbf{PONTIFÍCIA UNIVERSIDADE CATÓLICA DO PARANÁ}} \\[4pt]
  Escola Politécnica \\
  Curso: Ciência da Computação \\
  Disciplina: Métodos Quantitativos \\[10pt]
  {\large \textbf{Atividade A14 -- Teste de Hipótese}} \\[4pt]
  Aluno: JAFTE CARNEIRO FAGUNDES DA SILVA \\
  Data: 01/junho/2026
\end{center}

\vspace{0.5cm}
\hrule
\vspace{0.8cm}

% ═══════════════════════════════════════════════════════════════════════════════
\section*{Questão 1 -- Teste de Hipótese com Intervalo de Confiança}
% ═══════════════════════════════════════════════════════════════════════════════

\subsection*{Situação-Problema}

O gerente de TI de uma empresa de software afirma que o tempo médio de resposta
do sistema de atendimento ao cliente é de 3{,}0 segundos. Um analista de
qualidade suspeita que o tempo médio seja diferente disso. Para verificar, ele
coletou uma amostra aleatória de \textbf{45 requisições}, obtendo:

\begin{itemize}
  \item Média amostral: $\bar{X} = 3{,}42$ segundos
  \item Desvio padrão amostral: $S_X = 1{,}20$ segundos
\end{itemize}

Com \textbf{95\% de confiança}, o valor de 3{,}0 segundos pode ser considerado
representativo do tempo médio de resposta?

\subsection*{Solução}

\textbf{Passo 1: Hipóteses}

\[
  H_0: \mu = 3{,}0 \text{ s} \qquad H_1: \mu \neq 3{,}0 \text{ s}
\]

Trata-se de um teste \textbf{bilateral (duas caudas)}.

\textbf{Passo 2: Nível de significância}

\[
  \alpha = 0{,}05 \quad \Rightarrow \quad \text{valor crítico } Z = \pm 1{,}96
\]

\textbf{Passo 3: Intervalo de confiança}

Como $n = 45 > 30$, utiliza-se a distribuição $Z$ (normal):

\[
  \mu = \bar{X} \pm Z \cdot \frac{S_X}{\sqrt{n}}
      = 3{,}42 \pm 1{,}96 \cdot \frac{1{,}20}{\sqrt{45}}
      = 3{,}42 \pm 1{,}96 \cdot 0{,}1789
      = 3{,}42 \pm 0{,}3507
\]

\[
  \boxed{\mu \in \left(3{,}069\ ;\ 3{,}771\right)}
\]

\textbf{Passo 4: Decisão}

O valor $\mu_0 = 3{,}0$ \textbf{não está contido} no intervalo $(3{,}069\ ;\ 3{,}771)$.
Portanto, há evidência suficiente para \textbf{rejeitar $H_0$} ao nível de
significância de 5\%. O tempo médio de resposta é significativamente diferente
de 3{,}0 segundos.

\subsection*{Pseudocódigo}

\begin{algorithm}[H]
\caption{Teste de Hipótese -- Intervalo de Confiança (distribuição $Z$)}
\begin{algorithmic}[1]
  \State \textbf{Entrada:} $\bar{X}$, $S_X$, $n$, $\mu_0$, $\alpha$
  \State $Z_c \gets z\_critico(\alpha)$ \Comment{ex.: $\pm 1{,}96$ para $\alpha=0{,}05$}
  \State $erro \gets Z_c \times \dfrac{S_X}{\sqrt{n}}$
  \State $LI \gets \bar{X} - erro$
  \State $LS \gets \bar{X} + erro$
  \If{$\mu_0 < LI$ \textbf{ou} $\mu_0 > LS$}
    \State \textbf{Rejeitar} $H_0$
  \Else
    \State \textbf{Não rejeitar} $H_0$
  \EndIf
\end{algorithmic}
\end{algorithm}

\newpage

% ═══════════════════════════════════════════════════════════════════════════════
\section*{Questão 2 -- Diferença entre Médias: Amostras Grandes}
% ═══════════════════════════════════════════════════════════════════════════════

\subsection*{Situação-Problema}

Uma universidade quer verificar se existe diferença na quantidade média diária de
commits em dois repositórios de projetos distintos: o \textbf{Repositório A}
(projetos de graduação) e o \textbf{Repositório B} (projetos de pós-graduação).
Foram coletadas as seguintes informações:

\begin{center}
\begin{tabular}{lcc}
  \toprule
  & \textbf{Repositório A} & \textbf{Repositório B} \\
  \midrule
  Tamanho da amostra ($n$) & 50        & 60       \\
  Média ($\bar{X}$)        & 12{,}30   & 13{,}85  \\
  Variância ($S^2$)        & 8{,}00    & 12{,}50  \\
  \bottomrule
\end{tabular}
\end{center}

Com \textbf{significância de 5\%}, há diferença significativa entre as médias?

\subsection*{Solução}

\textbf{Hipóteses}

\[
  H_0: \mu_1 = \mu_2 \qquad H_1: \mu_1 \neq \mu_2
\]

\textbf{Estatística $Z_0$}

Como $n_1 = 50 > 30$ e $n_2 = 60 > 30$, e as variâncias populacionais são
desconhecidas, usa-se a aproximação pelas variâncias amostrais:

\[
  Z_0 = \frac{(\bar{X}_1 - \bar{X}_2) - (\mu_1 - \mu_2)}
             {\sqrt{\dfrac{S_1^2}{n_1} + \dfrac{S_2^2}{n_2}}}
      = \frac{(12{,}30 - 13{,}85) - 0}
             {\sqrt{\dfrac{8{,}00}{50} + \dfrac{12{,}50}{60}}}
\]

\[
  Z_0 = \frac{-1{,}55}{\sqrt{0{,}16 + 0{,}2083}}
      = \frac{-1{,}55}{\sqrt{0{,}3683}}
      = \frac{-1{,}55}{0{,}6069}
      \approx -2{,}55
\]

\textbf{Decisão}

O valor crítico para $\alpha = 0{,}05$ bilateral é $Z = \pm 1{,}96$.
Como $|{-2{,}55}| = 2{,}55 > 1{,}96$, o valor calculado está \textbf{fora}
da região de não-rejeição. Portanto, \textbf{rejeita-se $H_0$}: há diferença
significativa entre as médias de commits dos dois repositórios.

\subsection*{Pseudocódigo}

\begin{algorithm}[H]
\caption{Teste $Z$ -- Diferença entre Médias (Amostras Grandes)}
\begin{algorithmic}[1]
  \State \textbf{Entrada:} $\bar{X}_1, \bar{X}_2, S_1^2, S_2^2, n_1, n_2, \alpha$
  \State $erro\_padrao \gets \sqrt{S_1^2 / n_1 + S_2^2 / n_2}$
  \State $Z_0 \gets (\bar{X}_1 - \bar{X}_2) \;/\; erro\_padrao$
  \State $Z_c \gets z\_critico(\alpha)$ \Comment{ex.: $1{,}96$ para $\alpha=0{,}05$}
  \If{$|Z_0| > Z_c$}
    \State \textbf{Rejeitar} $H_0$
  \Else
    \State \textbf{Não rejeitar} $H_0$
  \EndIf
\end{algorithmic}
\end{algorithm}

\newpage

% ═══════════════════════════════════════════════════════════════════════════════
\section*{Questão 3 -- Amostras Pequenas e Variâncias Iguais}
% ═══════════════════════════════════════════════════════════════════════════════

\subsection*{Situação-Problema}

Um pesquisador deseja comparar o consumo médio de memória RAM (em MB) de dois
algoritmos de ordenação executados sob as mesmas condições. Assume-se que as
variâncias das populações são iguais. Os dados coletados foram:

\begin{center}
\begin{tabular}{lcc}
  \toprule
  & \textbf{Algoritmo A} & \textbf{Algoritmo B} \\
  \midrule
  Tamanho da amostra ($n$) & 10      & 12      \\
  Média ($\bar{X}$, em MB) & 128{,}4 & 134{,}7 \\
  Variância ($S^2$)        & 22{,}5  & 19{,}8  \\
  \bottomrule
\end{tabular}
\end{center}

Com \textbf{significância de 5\%}, há diferença no consumo médio de memória?

\subsection*{Solução}

\textbf{Hipóteses}

\[
  H_0: \mu_1 = \mu_2 \qquad H_1: \mu_1 \neq \mu_2
\]

\textbf{Variância agrupada} $S_p^2$

\[
  S_p^2 = \frac{(n_1 - 1)\,S_1^2 + (n_2 - 1)\,S_2^2}{n_1 + n_2 - 2}
        = \frac{(9)(22{,}5) + (11)(19{,}8)}{10 + 12 - 2}
        = \frac{202{,}5 + 217{,}8}{20}
        = \frac{420{,}3}{20}
        = 21{,}015
\]

\textbf{Estatística $t_0$}

\[
  t_0 = \frac{(\bar{X}_1 - \bar{X}_2) - (\mu_1 - \mu_2)}
             {\sqrt{S_p^2 \left(\dfrac{1}{n_1} + \dfrac{1}{n_2}\right)}}
      = \frac{(128{,}4 - 134{,}7) - 0}
             {\sqrt{21{,}015 \left(\dfrac{1}{10} + \dfrac{1}{12}\right)}}
\]

\[
  = \frac{-6{,}3}{\sqrt{21{,}015 \times 0{,}1833}}
  = \frac{-6{,}3}{\sqrt{3{,}852}}
  = \frac{-6{,}3}{1{,}9627}
  \approx -3{,}21
\]

\textbf{Graus de liberdade:} $gl = n_1 + n_2 - 2 = 10 + 12 - 2 = 20$

Para $\alpha = 0{,}05$ bilateral e $gl = 20$, o valor crítico é
$t_{(20;\,0{,}025)} \approx 2{,}086$.

\textbf{Decisão}

Como $|{-3{,}21}| = 3{,}21 > 2{,}086$, \textbf{rejeita-se $H_0$}: há diferença
significativa no consumo médio de memória entre os dois algoritmos.

\subsection*{Pseudocódigo}

\begin{algorithm}[H]
\caption{Teste $t$ -- Amostras Pequenas, Variâncias Iguais (variância agrupada)}
\begin{algorithmic}[1]
  \State \textbf{Entrada:} $\bar{X}_1, \bar{X}_2, S_1^2, S_2^2, n_1, n_2, \alpha$
  \State $S_p^2 \gets \dfrac{(n_1-1)\,S_1^2 + (n_2-1)\,S_2^2}{n_1 + n_2 - 2}$
  \State $t_0 \gets \dfrac{\bar{X}_1 - \bar{X}_2}{\sqrt{S_p^2\,(1/n_1 + 1/n_2)}}$
  \State $gl \gets n_1 + n_2 - 2$
  \State $t_c \gets t\_critico(\alpha,\; gl)$
  \If{$|t_0| > t_c$}
    \State \textbf{Rejeitar} $H_0$
  \Else
    \State \textbf{Não rejeitar} $H_0$
  \EndIf
\end{algorithmic}
\end{algorithm}

\newpage

% ═══════════════════════════════════════════════════════════════════════════════
\section*{Questão 4 -- Amostras Pequenas e Variâncias Diferentes}
% ═══════════════════════════════════════════════════════════════════════════════

\subsection*{Situação-Problema}

Dois grupos de dispositivos IoT de fabricantes distintos foram testados quanto
ao tempo médio de resposta (em ms) a uma requisição MQTT. As variâncias
populacionais são presumivelmente \textbf{diferentes}. Os dados foram:

\begin{center}
\begin{tabular}{lcc}
  \toprule
  & \textbf{Fabricante X} & \textbf{Fabricante Y} \\
  \midrule
  Tamanho da amostra ($n$) & 8       & 10      \\
  Média ($\bar{X}$, em ms) & 42{,}5  & 48{,}3  \\
  Variância ($S^2$)        & 18{,}0  & 62{,}4  \\
  \bottomrule
\end{tabular}
\end{center}

Com \textbf{significância de 5\%}, os tempos médios de resposta diferem?

\subsection*{Solução}

\textbf{Hipóteses}

\[
  H_0: \mu_1 = \mu_2 \qquad H_1: \mu_1 \neq \mu_2
\]

\textbf{Estatística $t^*$}

\[
  t^* = \frac{\bar{X}_1 - \bar{X}_2}
             {\sqrt{\dfrac{S_1^2}{n_1} + \dfrac{S_2^2}{n_2}}}
      = \frac{42{,}5 - 48{,}3}
             {\sqrt{\dfrac{18{,}0}{8} + \dfrac{62{,}4}{10}}}
      = \frac{-5{,}8}{\sqrt{2{,}25 + 6{,}24}}
      = \frac{-5{,}8}{\sqrt{8{,}49}}
      = \frac{-5{,}8}{2{,}9138}
      \approx -1{,}99
\]

\textbf{Graus de liberdade ajustados} (Welch--Satterthwaite):

\[
  gl = \frac{\left(\dfrac{S_1^2}{n_1} + \dfrac{S_2^2}{n_2}\right)^2}
            {\dfrac{\left(\dfrac{S_1^2}{n_1}\right)^2}{n_1 - 1}
             + \dfrac{\left(\dfrac{S_2^2}{n_2}\right)^2}{n_2 - 1}}
     = \frac{(2{,}25 + 6{,}24)^2}
            {\dfrac{(2{,}25)^2}{7} + \dfrac{(6{,}24)^2}{9}}
     = \frac{(8{,}49)^2}{\dfrac{5{,}0625}{7} + \dfrac{38{,}9376}{9}}
\]

\[
  = \frac{72{,}0801}{0{,}7232 + 4{,}3264}
  = \frac{72{,}0801}{5{,}0496}
  \approx 14{,}27
  \quad \Rightarrow \quad gl = 14 \text{ (truncado)}
\]

Para $\alpha = 0{,}05$ bilateral e $gl = 14$, o valor crítico é
$t_{(14;\,0{,}025)} \approx 2{,}145$.

\textbf{Decisão}

Como $|{-1{,}99}| = 1{,}99 < 2{,}145$, o valor calculado está \textbf{dentro}
da região de não-rejeição. Portanto, \textbf{não se rejeita $H_0$}: não há
evidência suficiente para afirmar que os tempos médios de resposta dos dois
fabricantes são diferentes, ao nível de significância de 5\%.

\subsection*{Pseudocódigo}

\begin{algorithm}[H]
\caption{Teste $t^*$ -- Amostras Pequenas, Variâncias Diferentes (Welch)}
\begin{algorithmic}[1]
  \State \textbf{Entrada:} $\bar{X}_1, \bar{X}_2, S_1^2, S_2^2, n_1, n_2, \alpha$
  \State $A \gets S_1^2 / n_1$
  \State $B \gets S_2^2 / n_2$
  \State $t^* \gets (\bar{X}_1 - \bar{X}_2) \;/\; \sqrt{A + B}$
  \State $gl \gets \lfloor (A + B)^2 \;/\; (A^2/(n_1-1) + B^2/(n_2-1)) \rfloor$
  \State $t_c \gets t\_critico(\alpha,\; gl)$
  \If{$|t^*| > t_c$}
    \State \textbf{Rejeitar} $H_0$
  \Else
    \State \textbf{Não rejeitar} $H_0$
  \EndIf
\end{algorithmic}
\end{algorithm}

\end{document}
